\documentclass[a4paper, 14pt]{extarticle}

\usepackage{cmap}
\usepackage[T2A]{fontenc}
\usepackage[utf8]{inputenc}
\usepackage[english, russian]{babel}
\usepackage{amsmath}
\usepackage{subcaption}
\usepackage{graphicx}
\usepackage{float}
\usepackage{hyperref}
\usepackage{multirow}

\usepackage[left=3cm, right=1.5cm, top=2cm, bottom=2cm]{geometry}
\usepackage{setspace}
\onehalfspacing 

\usepackage{indentfirst}
\setlength{\parindent}{1.25cm} 
\setlength{\emergencystretch}{3em}
\usepackage[expansion=false,protrusion=false,nopatch=footnote]{microtype}

\usepackage{tocloft}
\addto\captionsrussian{\renewcommand{\contentsname}{СОДЕРЖАНИЕ}}
\renewcommand{\cfttoctitlefont}{\hfill\normalfont}
\renewcommand{\cftaftertoctitle}{\hfill\vspace{\baselineskip}}
\renewcommand{\cftsecdotsep}{\cftdotsep}

\begin{document}

\begin{titlepage}
    \begin{center}
        Министерство науки и высшего образования Российской Федерации\\
        Федеральное государственное автономное образовательное учреждение высшего образования\\
        «Санкт-Петербургский политехнический университет Петра Великого»\\
        
        \textbf{Физико-механический институт}\\
        \textbf{Высшая школа прикладной математики и вычислительной физики}
    \end{center}

    \vfill
    
    \begin{center}
        \textbf{Отчёт по лабораторной работе \textnumero~ 6\\[0.5cm]}
        \textbf{Линейная регрессия. Метод наименьших квадратов и наименьших модулей\\}
        по дисциплине <<Математическая статистика>>
    \end{center}

    \vfill

    \hfill
    \begin{minipage}{0.45\textwidth}
        \textbf{Выполнил студент гр. 5030102/30202:}\\
        Шильдт Д.С.\\[0.5cm]
        \textbf{Преподаватель:}\\ 
        Баженов А.Н.
    \end{minipage}

    \vfill

    \begin{center}
        Санкт-Петербург\\2026
    \end{center}
\end{titlepage}

\setcounter{page}{2}

\tableofcontents
\newpage

\section{Постановка задачи}

\textbf{Задание:}
\begin{enumerate}
    \item Задать равномерные значения $x_i$ на отрезке $[-1.8, 2]$ с шагом 0.2 (20 точек).
    \item Сгенерировать $y_i= 2 + 2x_i + \varepsilon_i$, где $\varepsilon_i \sim N(0, 1)$.
    \item Оценить параметры $a,b$ двумя методами:
    \begin{itemize}
        \item  Метод наименьших квадратов (МНК).
        \item  Метод наименьших модулей (МНМ).
    \end{itemize}
    \item Повторить пункты 1-3, добавив выбросы: $y_1\leftarrow y_1 + 10,\ y_{20}\leftarrow y_{20} - 10$
\end{enumerate}

\textbf{Цели и задачи:}
\begin{itemize}
    \item Сравнить устойчивость МНК и МНМ к выбросам.
    \item  Построить графики регрессионных прямых и исходных данных.
\end{itemize}

\section{Теоретическая часть}

Регрессионную модель описания данных называют \textit{простой линейной регрессией}, если
\begin{equation*}
    y_i = \beta_0 + \beta_1 x_i + \varepsilon_i, \quad i = 1, \dots, n,
\end{equation*}
где $x_1,\ ...\ ,\ x_n$ --- заданные числа (значения фактора); $y_1,\ ...\ ,\ y_n$ --- наблюдаемые значения отклика; $\varepsilon_1,\ ...\ ,\ \varepsilon_n$ --- независимые, нормально распределённые $N(0,\sigma)$ с нулевым математическим ожиданием и одинаковой (неизвестной) дисперсией случайные величины (ненаблюдаемые); $\beta_0,\beta_1$ --- неизвестные параметры, подлежащие оцениванию.


Метод наименьших квадратов (МНК):
\begin{equation*}
    Q(\beta_0, \beta_1) = \sum_{i=1}^n \varepsilon^2_i = \sum_{i=1}^n (y_i - \beta_0 - \beta_1 x_i)^2 \to \min_{\beta_0, \beta_1}.
\end{equation*}

МНК-оценки параметров $\widehat{\beta_0}$ и $\widehat{\beta_1}$ находятся из условия обращения функции $Q(\beta_0,\beta_1)$ в минимум.
\begin{align*}
    \widehat{\beta}_1 &= \frac{\overline{xy} - \bar{x} \cdot \bar{y}}{\overline{x^2} - (\bar{x})^2}, & \widehat{\beta}_0 &= \bar{y} - \bar{x}\widehat{\beta}_1.
\end{align*}

Метод наименьших модулей (МНМ):
\begin{equation*}
    M(\beta_0, \beta_1) = \sum_{i=1}^n |y_i - \beta_0 - \beta_1 x_i| \to \min_{\beta_0, \beta_1}.
\end{equation*}

\section{Реализация}

\subsection{Язык программирования и среда}
\begin{itemize}
    \item \textbf{Используемый язык:} Python 3.14.2
    \item \textbf{Используемая среда разработки:} Visual Studio Code
\end{itemize}

\subsection{Используемые библиотеки}
\begin{itemize}
    \item \textbf{numpy} --- применяется для генерации шума $\varepsilon_i$, векторных вычислений и расчета сумм квадратов и смешанных произведений при аналитическом решении МНК.
    \item \textbf{scipy.optimize} --- используется (функция \texttt{minimize}) для численной минимизации негладкой целевой функции $L_1$-нормы в методе наименьших модулей.
    \item \textbf{matplotlib.pyplot} --- обеспечивает визуализацию исходных данных, аномальных наблюдений и расчетных регрессионных прямых.
    \item \textbf{pathlib} --- определение путей для экспорта числовых метрик и графиков.
\end{itemize}

\subsection{Суть реализованных методов}

\begin{itemize}
    \item \texttt{generate\_x\_values()} и \texttt{generate\_response(x, rng)} --- формируют равномерную сетку детерминированного фактора $x_i \in [-1.8, 2.0]$ и вектор по модели $y_i = 2 + 2x_i + \varepsilon_i$.
    \item \texttt{add\_outliers(y)} --- искусственно модифицирует крайние порядковые элементы массива, прибавляя $10$ к первому значению и вычитая $10$ из последнего ($y_1 \leftarrow y_1 + 10$, $y_{20} \leftarrow y_{20} - 10$).
    \item \texttt{fit\_ols(x, y)} --- реализует метод наименьших квадратов.
    \item \texttt{fit\_lad(x, y)} --- реализует метод наименьших модулей.
    \item \texttt{parameter\_errors(estimate, true\_value)} --- модуль, вычисляющий отклонение и относительную погрешность.
    \item \texttt{evaluate\_methods()} и \texttt{save\_regression\_figures(...)} --- производят расчеты для двух выборок, с последующим сохранением графиков и данных в (CSV).
\end{itemize}

\section{Результаты}

Вычисленные оценки параметров сдвига $\hat{a}$ и масштаба $\hat{b}$, а также соответствующие значения абсолютной $\Delta z$ и относительной $\delta z$ погрешностей для двух экспериментов сведены в таблицу \ref{tab:reg_results}. Теоретические значения коэффициентов генеральной совокупности строго заданы: $a = 2$, $b = 2$.

\begin{table}[H]
    \centering
    \caption{Оценки параметров линейной регрессии и их погрешности}
    \label{tab:reg_results}
    \begin{tabular}{|l|l|c|c|c|c|c|c|}
        \hline
        \textbf{Сценарий} & \textbf{Метод} & $\mathbf{\hat{a}}$ & $\mathbf{\Delta a}$ & $\mathbf{\delta a, \%}$ & $\mathbf{\hat{b}}$ & $\mathbf{\Delta b}$ & $\mathbf{\delta b, \%}$ \\ \hline
        \multirow{2}{*}{Без выбросов} 
        & МНК & $1.9544$ & $0.0456$ & $2.28$ & $2.1269$ & $0.1269$ & $6.35$ \\ \cline{2-8}
        & МНМ & $2.0969$ & $0.0969$ & $4.85$ & $1.9485$ & $0.0515$ & $2.57$ \\ \hline
        \multirow{2}{*}{С выбросами} 
        & МНК & $2.0972$ & $0.0972$ & $4.86$ & $0.6984$ & $1.3016$ & $65.08$ \\ \cline{2-8}
        & МНМ & $2.0969$ & $0.0969$ & $4.85$ & $1.9485$ & $0.0515$ & $2.58$ \\ \hline
    \end{tabular}
\end{table}

На выборке без выбросов (рис. \ref{fig:reg_clean}) методы МНК и МНМ формируют прямые, визуально и аналитически близкие к теоретической модели. При внесении выбросов амплитудой $\pm 10$ в крайние точки $y_1$ и $y_{20}$ (рис. \ref{fig:reg_outliers}) прямая МНК меняет угол наклона, пытаясь минимизировать возросшую сумму квадратов остатков. Это количественно подтверждается скачком относительной погрешности параметра масштаба $\delta b$ с $6.35\%$ до $65.08\%$. Прямая МНМ, опирающаяся на оптимизацию суммы модулей, полностью игнорирует амплитуду аномалий и сохраняет свое положение в пространстве неизменным ($\hat{b} \approx 1.9485$ в обоих сценариях).

\begin{figure}[H]
    \centering
    \includegraphics[width=0.8\textwidth]{../figures/regression_clean.png}
    \caption{Линейная регрессия для выборки без выбросов}
    \label{fig:reg_clean}
\end{figure}

\begin{figure}[H]
    \centering
    \includegraphics[width=0.8\textwidth]{../figures/regression_outliers.png}
    \caption{Линейная регрессия для выборки с выбросами}
    \label{fig:reg_outliers}
\end{figure}

\section{Обсуждение}

При анализе выборки, не содержащей выбросов, оба алгоритма показывают высокую точность восстановления истинной зависимости. Относительные погрешности вычислений не превышают $6.35\%$. Метод наименьших квадратов (МНК) дает более точную оценку параметра сдвига ($\delta a = 2.28\%$ против $4.85\%$ у МНМ), тогда как метод наименьших модулей (МНМ) точнее оценивает масштаб ($\delta b = 2.57\%$ против $6.35\%$). Полученные значения подтверждают сопоставимую эффективность оптимизации $L_2$- и $L_1$-норм при работе со стандартным гауссовским шумом без грубых искажений [1].

Внесение выбросов амплитудой $\pm 10$ меняет поведение алгоритмов. Квадратичная функция потерь МНК придает аномальным отклонениям доминирующий вес. Это заставляет регрессионную прямую изменить угол наклона для минимизации суммы квадратов остатков на краях отрезка. Оценка параметра масштаба $\hat{b}$ падает с $2.1269$ до $0.6984$, а относительная погрешность достигает $65.08\%$. Метод наименьших квадратов демонстрирует абсолютную неробастность и непригодность для анализа зашумленных данных.

Метод наименьших модулей выявляет строгую устойчивость к выбросам. Переход к метрике $L_1$ ограничивает влияние выбросов: итеративная целевая функция реагирует на направление отклонения, полностью игнорируя его амплитуду. Оценки параметров МНМ на выборке с выбросами остаются практически идентичными результатам чистого эксперимента ($\hat{a} \approx 2.0969$, $\hat{b} \approx 1.9485$). Метод локализует основную линейную структуру данных, полностью сглаживая искажающий эффект грубых ошибок.

\section{Выводы}

В ходе выполнения работы установлено, что на выборке без выбросов методы наименьших квадратов (МНК) и наименьших модулей (МНМ) обеспечивают высокую и сопоставимую точность восстановления параметров линейной регрессии. Максимальная относительная погрешность вычислений состаила $6.35\%$ для оценки параметра масштаба.

Метод наименьших квадратов обладает абсолютной неробастностью [1]. Внесение выбросов амплитудой $\pm 10$ приводит к искажению аппроксимирующей прямой. Из-за минимизации $L_2$-нормы (суммы квадратов остатков) алгоритм придает аномалиям доминирующий вес: оценка наклона падает до $\hat{b} = 0.6984$, а относительная погрешность достигает $\delta b = 65.08\%$.

Метод наименьших модулей демонстрирует строгую устойчивость к грубым искажениям отклика [1]. Оптимизация на основе $L_1$-нормы (суммы модулей остатков) позволяет полностью сгладить влияние амплитуды аномалий на положение регрессионной прямой. На зашумленной выборке оценка масштаба составляет $\hat{b} \approx 1.9485$ при погрешности $\delta b = 2.58\%$, что практически схоже с результатами чистого эксперимента ($\delta b = 2.57\%$).
 
\refstepcounter{section}
\addcontentsline{toc}{section}{\thesection\ Список литературы}
\renewcommand{\refname}{\thesection\ Список литературы}

\begin{thebibliography}{9}

    \bibitem{babakhina}
    Бабахина С. А., Баженов А. Н. Практикум по математической статистике: учеб. пособие. — СПб. : Санкт-Петербургский политехнический университет Петра Великого, 2026. 

    \bibitem{maksimov}
    Максимов Ю. Д. Математика. Теория и практика по математической статистике. Конспект-справочник по теории вероятностей : учеб. пособие. — СПб. : Изд-во Политехн. ун-та, 2009. — 395 с. 

\end{thebibliography}

\section{Приложение}

Исходный код программ доступен в репозитории на GitHub: \\
\url{https://github.com/Reichenau/math-stats-labs}
 
\end{document}
